\section{A defect in the staged fitting schedule}
\label{sec:staged}

The registration stage (\S\ref{sec:registration}) runs as four sequential
stages (initialisation, posing/transforming, coarse vertex deform, fine
vertex deform), and the audit around \F{F4} and \F{F6} also exposed
a defect in how this schedule allocates iterations: pose was held fixed for
2000 of the refinement stage's 2400 iterations, with only the first 400
free to adjust joint rotation before vertex deformation took over. Because
Chamfer distance descends smoothly regardless of when pose freezes relative
to deformation, this defect produced no visible signature in the loss curve
and survived undetected until this investigation. A rebuilt schedule
widening the pose-adjustable window was tested against the stock schedule on
the full integrity-metric suite of \S\ref{sec:eval-suite}: it improves
nearly every metric measured, at no additional computational cost, and was
shipped into the delivered recipe. This is a second, independent instance of
the pattern in \F{F1}: the objective's own value is not a reliable signal
for whether the fitting procedure that produced it was sound, caught only
because integrity metrics, not surface distance alone, were checked.
